\subsection{Branchwise strict optimization}

Because $F$ has different algebraic expressions on different contact cells
and has a genuine jump at one high-radial boundary, formalization must not
silently optimize one expression across another expression's domain.

\begin{definition}[Strict branch certificate]
\label{def:s2-pure-strict-certificate}
Let $D$ be a bounded domain, let
\[
 D=\bigcup_{\lambda\in\Lambda}D_\lambda
\]
be a finite disjoint branch decomposition, and let
$\Psi_\lambda$ be the algebraic-radical expression for the objective on
$D_\lambda$.  A \emph{strict branch certificate} for $\Psi<0$ consists, for
every $\lambda\in\Lambda$, of
\[
 \max_{z\in\overline{D_\lambda}}\Psi_\lambda(z)\le0
\]
together with
\[
 D_\lambda\cap\{z:\Psi_\lambda(z)=0\}=\varnothing.
\]
The second condition is essential.  On an open feasible domain the supremum
may equal zero even when the desired pointwise inequality is strict, so a
claim of a uniform positive margin must not be inserted unless it is proved
separately.
\end{definition}

For Problems~\ref{prob:s2-one-gap-endpoint}--\ref{prob:s2-ce2-five-chain},
the branch index records the selected case of every occurrence of $F$.
Squaring a radical inequality is allowed only after its sign hypotheses have
been added to the same branch domain.

For a polynomial-only encoding, replace each radical $\sqrt{P(z)}$ by a fresh
variable $s$ with
\[
 s\ge0,\qquad s^2=P(z).
\]
Replace every minimum by its two order branches, retaining equality in one
specified branch only.  All denominators displayed in
Definition~\ref{def:s2-pure-map} are positive on their selected domains; that
positivity should be recorded before clearing denominators.  These steps
produce a finite semialgebraic problem without changing the feasible
component.

\subsection{One signed parameter domain}

Use three independent variables
\[
 (r,a,d)\in\mathbb R^3
\]
and define
\[
 w=1-r,\qquad
 e=\sqrt{1-rw},\qquad
 \eta=1-e,\qquad
 p=e(1-e),\qquad
 k=\eta+a+d,
\]
\[
 \Delta_R=p-a-wd,
 \qquad
 \Delta_L=p-ra-d.
\]

\begin{definition}[Pure signed domains]
\label{def:s2-pure-signed-domains}
Define
\[
 \mathcal D_{\rm sgn}
 =
 \left\{
 (r,a,d):
 0<r<1,\ a>0,\ d>0,\ \Delta_R>0
 \right\}.
\]
Define
\[
 \mathcal D_1
 =
 \left\{
 (r,a,d)\in\mathcal D_{\rm sgn}:
 \Delta_L\le0,\quad
 a<\frac w2,\quad
 d<\frac r2,\quad
 rd-wa>0
 \right\},
\]
and
\[
 \mathcal D_2
 =
 \left\{
 (r,a,d)\in\mathcal D_{\rm sgn}:
 \Delta_L>0
 \right\}.
\]
\end{definition}

\begin{proposition}[Well-definedness on the pure domains]
\label{prop:s2-pure-domain-bounds}
Every occurrence of $F$ in
Problems~\ref{prob:s2-one-gap-endpoint}--\ref{prob:s2-ce2-five-chain} has first argument in $(1/2,1)$ and second
argument in $[0,1]$.
\end{proposition}

\begin{proof}
Since
\[
 1-e=\frac{rw}{1+e},
 \qquad
 p=\frac{erw}{1+e}<\frac{rw}{2},
\]
the inequality $\Delta_R>0$ gives
\[
 0<a<p<\frac{rw}{2},
 \qquad
 0<d<\frac p w<\frac r2.
\]
Consequently
\[
 1-\frac a w>\frac12,\qquad
 1-\frac d r>\frac12,\qquad
 \frac{w-a}{w}>\frac12,\qquad
 \frac{r-d}{r}>\frac12.
\]
Also
\[
 w+d-\frac{k}{r}=\frac{\Delta_R}{r}>0,
\]
so
\[
 0<\frac{k}{r}<w+d<1.
\]
Thus both endpoint problems are well-defined.  For the return problems,
\[
 0<a<\frac12,\qquad
 0<d<\frac12,\qquad
 0<m=\min\left\{\frac a r,\frac d w\right\}<\frac12,
\]
and
\[
 0<Y_0=\frac{k}{2r}<\frac12.
\]
The update $T$ maps into $[0,1]$, so the remaining input bounds follow
inductively.
\end{proof}

The bridge from the geometric notation is
\[
 r=R,\qquad a=\alpha,\qquad d=\delta.
\]
This line is the entire geometry-to-analysis dictionary for the four
formalization-ready problems.

